\chapter{Theoretical Background}

% ─────────────────────────────────────────────────────────────────────────────
\section{SLAM}

SLAM (Simultaneous Localisation and Mapping) estimates the robot trajectory
$\mathbf{x}_{1:t}$ and the map $\mathbf{m}$ jointly from sensor measurements
$\mathbf{z}_{1:t}$ and motion commands $\mathbf{u}_{1:t}$~\cite[p.~245]{thrun2005probabilistic}:

\begin{equation}
  p(\mathbf{x}_{1:t},\, \mathbf{m} \mid \mathbf{z}_{1:t},\, \mathbf{u}_{1:t}).
  \label{eq:slam_full}
\end{equation}

\noindent Filter-based methods solve this by maintaining a state distribution
updated at each measurement; graph-based methods retain the full trajectory and
enforce global consistency at loop closure.
These two families are described in Sections~\ref{sec:filters}
and~\ref{sec:graph_slam}.

% ─────────────────────────────────────────────────────────────────────────────
\section{Sensors}
\label{sec:sensors}

This thesis uses four sensor types: a 3D LiDAR, a stereo camera, an IMU, and
wheel odometry encoders.

\subsection{LiDAR}

A LiDAR measures distances by emitting laser pulses and recording the
round-trip travel time; the range is $d = c\,\Delta t / 2$.
A rotating 3D model fires pulses across multiple vertical channels while
spinning horizontally, producing a cloud of tens of thousands of $(x,y,z)$
points per revolution~\cite{siegwart2004introduction}.
Its advantages are metric accuracy and independence from ambient lighting; the
main limitation is motion distortion, as points within one scan are captured at
slightly different robot poses.

\subsection{Stereo Camera}

A stereo camera recovers depth from the horizontal disparity $d$ between
projections of a scene point in two calibrated images separated by baseline $B$:

\begin{equation}
  Z = \frac{f \cdot B}{d},
  \label{eq:stereo_depth}
\end{equation}

\noindent where $f$ is the focal length in pixels~\cite{siegwart2004introduction}.
Cameras provide rich texture for feature matching but performance degrades in
low-light or textureless environments.

\subsection{IMU}

An IMU combines a three-axis accelerometer and gyroscope whose outputs are
corrupted by slowly varying biases and noise~\cite{thrun2005probabilistic}:

\begin{equation}
  \tilde{\mathbf{a}} = \mathbf{a} + \mathbf{b}_a + \boldsymbol{\eta}_a, \qquad
  \tilde{\boldsymbol{\omega}} = \boldsymbol{\omega} + \mathbf{b}_g + \boldsymbol{\eta}_g.
  \label{eq:imu_model}
\end{equation}

\noindent The IMU runs at high frequency, providing dense motion information
between slower sensor updates, but must be fused with another sensor to prevent
drift from accumulating.

\subsection{Wheel Odometry}

Wheel encoders estimate the pose increment from the angular velocities of each
wheel. For a differential-drive robot at heading $\theta$~\cite{siegwart2004introduction}:

\begin{equation}
  \Delta x = \frac{v_L + v_R}{2}\,\Delta t\cos\theta, \quad
  \Delta y = \frac{v_L + v_R}{2}\,\Delta t\sin\theta, \quad
  \Delta\theta = \frac{v_R - v_L}{L}\,\Delta t,
  \label{eq:wheel_odom}
\end{equation}

\noindent where $L$ is the wheel separation. Odometry is a useful short-term
prior but drifts due to wheel slip and calibration errors.

% ─────────────────────────────────────────────────────────────────────────────
\section{Map Representations}
\label{sec:map_representations}

\subsection{Occupancy Grid}

An occupancy grid divides the 2D plane into cells, each updated via the
log-odds of occupancy~\cite{thrun2005probabilistic}:

\begin{equation}
  l_{t,i} = \log \frac{p(m_i \mid \mathbf{z}_{1:t},\, \mathbf{x}_{1:t})}
                      {1 - p(m_i \mid \mathbf{z}_{1:t},\, \mathbf{x}_{1:t})}.
  \label{eq:log_odds}
\end{equation}

\noindent Cells above a threshold are marked occupied, those below are free,
and the remainder stay unknown.
Occupancy grids directly support path planning but are limited to 2D with
fixed cell resolution.

\subsection{3D Point Cloud}

A 3D point cloud accumulates LiDAR scan points in a common reference frame,
capturing the full three-dimensional geometry of the
environment~\cite{thrun2005probabilistic}.
It supports geometric evaluation against a reference scan but has higher memory
usage than an occupancy grid and contains no free-space information.
Visual SLAM systems produce sparse point clouds of triangulated image features;
these are not evaluated geometrically in this thesis.

% ─────────────────────────────────────────────────────────────────────────────
\section{Point Cloud Registration}

Point cloud registration finds the rigid transformation that best aligns two
overlapping point clouds; in LiDAR SLAM each new scan is aligned to the
previous one or to the accumulated map to estimate motion.

\subsection{LiDAR Feature Extraction}

Many algorithms reduce computation by extracting geometrically distinctive
features before registration~\cite{zhang2014loam}.
The local smoothness of point $i$ is estimated from its neighbours:

\begin{equation}
  c_i = \frac{1}{\lvert \mathcal{S}_i \rvert \cdot \lVert \mathbf{p}_i \rVert}
        \left\lVert \sum_{j \in \mathcal{S}_i} (\mathbf{p}_j - \mathbf{p}_i) \right\rVert.
  \label{eq:smoothness}
\end{equation}

\noindent Points with high curvature become edge features; those with low
curvature become planar features. FAST-LIO2 and LIO-SAM both use this step.

\subsection{Iterative Closest Point (ICP)}

ICP~\cite{besl1992icp} alternates between matching each source point to its
nearest neighbour and minimising:

\begin{equation}
  E(\mathbf{R}, \mathbf{t}) = \sum_i \lVert \mathbf{R}\mathbf{p}_i + \mathbf{t} - \mathbf{q}_i \rVert^2,
  \label{eq:icp}
\end{equation}

\noindent solved in closed form via SVD. ICP is sensitive to initialisation:
a poor starting estimate can cause convergence to a local minimum.

\subsection{Point-to-Plane ICP}

Point-to-plane ICP~\cite{chen1991point} minimises the distance from each
source point to the tangent plane at its correspondence:

\begin{equation}
  E(\mathbf{R}, \mathbf{t}) =
    \sum_i \bigl((\mathbf{R}\mathbf{p}_i + \mathbf{t} - \mathbf{q}_i)
                 \cdot \hat{\mathbf{n}}_i \bigr)^2.
  \label{eq:icp_plane}
\end{equation}

\noindent Sliding along surfaces incurs no penalty, so convergence is faster
than point-to-point ICP on planar regions.

\subsection{Generalized ICP (GICP)}

GICP~\cite{segal2009gicp} models each point as a Gaussian over its local
surface geometry and minimises the sum of Mahalanobis distances between
corresponding distributions, subsuming both ICP variants.
It is more robust to sensor noise and is the registration backend of KISS-ICP.

\subsection{Normal Distributions Transform (NDT)}

NDT~\cite{biber2003ndt} subdivides the target into voxels, each summarised by a
Gaussian $(\boldsymbol{\mu}_k, \boldsymbol{\Sigma}_k)$, and maximises:

\begin{equation}
  \text{score}(\mathbf{R}, \mathbf{t}) =
    \sum_i \exp\!\left(
      -\tfrac{1}{2}(\mathbf{R}\mathbf{p}_i + \mathbf{t} - \boldsymbol{\mu}_k)^\top
      \boldsymbol{\Sigma}_k^{-1}
      (\mathbf{R}\mathbf{p}_i + \mathbf{t} - \boldsymbol{\mu}_k)
    \right).
  \label{eq:ndt}
\end{equation}

\noindent The fixed voxel structure gives predictable memory usage and robustness
to non-uniform point density. MOLA Odometry uses NDT.

% ─────────────────────────────────────────────────────────────────────────────
\section{Filter-Based State Estimation}
\label{sec:filters}

Filter-based methods maintain a probability distribution over the robot state
and update it recursively via predict and correct steps.

\subsection{Extended Kalman Filter (EKF)}

The EKF~\cite{thrun2005probabilistic} represents the state as a Gaussian
$(\boldsymbol{\mu}_t, \boldsymbol{\Sigma}_t)$ and linearises nonlinear models
at the current estimate via Jacobians $\mathbf{F}_t$ and $\mathbf{H}_t$.
The prediction step is:

\begin{equation}
  \bar{\boldsymbol{\mu}}_t = f(\boldsymbol{\mu}_{t-1}, \mathbf{u}_t), \qquad
  \bar{\boldsymbol{\Sigma}}_t = \mathbf{F}_t \boldsymbol{\Sigma}_{t-1} \mathbf{F}_t^\top + \mathbf{Q}_t.
  \label{eq:ekf_predict}
\end{equation}

\noindent The correction step incorporates the measurement through the Kalman
gain. The \texttt{robot\_localization} package uses the EKF for dead-reckoning
in this thesis.

\subsection{Iterated Extended Kalman Filter (iEKF)}

The iEKF re-linearises the observation model around the updated estimate and
repeats the correction step until convergence, typically within 2 to 5
iterations, improving accuracy for nonlinear systems.
FAST-LIO2~\cite{xu2022fastlio2} uses an iEKF on the $SO(3)$ manifold to correct
pose, velocity, and IMU biases jointly.

\subsection{Rao-Blackwellized Particle Filter}

The RBPF~\cite{thrun2005probabilistic} factors the SLAM posterior by sampling
robot poses while computing each particle's map analytically.
Particles are weighted by:

\begin{equation}
  w_t^{[k]} \propto p\!\left(\mathbf{z}_t \mid \mathbf{x}_t^{[k]},\, \mathbf{m}^{[k]}\right),
  \label{eq:rbpf_weight}
\end{equation}

\noindent and low-weight particles are resampled. Loop closure emerges implicitly
as weights shift when the robot revisits a known area. GMapping uses this approach.

% ─────────────────────────────────────────────────────────────────────────────
\section{Graph-Based SLAM}
\label{sec:graph_slam}

Graph-based SLAM retains the full trajectory and enforces global consistency
whenever a loop closure is detected, unlike filter-based methods which discard
past information.

\subsection{Pose Graphs}

A pose graph~\cite{grisetti2010tutorial} stores a pose $\mathbf{x}_i \in SE(3)$
at each node; edges encode relative constraints and their information matrices
$\boldsymbol{\Omega}_{ij}$. Optimisation minimises:

\begin{equation}
  F(\mathbf{x}) = \sum_{(i,j)\in\mathcal{E}}
    \mathbf{e}_{ij}^\top\, \boldsymbol{\Omega}_{ij}\, \mathbf{e}_{ij},
  \label{eq:pose_graph}
\end{equation}

\noindent using sparse solvers such as g2o or GTSAM. Cartographer 2D, SLAM
Toolbox, and RTABMap all use pose graph back-ends.

\subsection{Factor Graphs}

A factor graph~\cite{cadena2016past} generalises pose graphs by allowing nodes
to represent any variable (poses, IMU biases, landmarks) and edges to encode
arbitrary probability distributions.
This is the natural framework for tightly coupled sensor fusion.
LIO-SAM uses a factor graph combining IMU preintegration, LiDAR odometry, loop
closure, and GPS constraints.

\subsection{Loop Closure Detection}

Loop closure detection~\cite{cadena2016past} adds long-range constraints when
the robot revisits a known location, correcting drift that otherwise grows
indefinitely.
For LiDAR systems, candidate pairs are verified by ICP; for visual systems, a
bag-of-words index enables place recognition across the full map history
(Section~\ref{sec:bow}).

% ─────────────────────────────────────────────────────────────────────────────
\section{IMU Pre-Integration and Deskewing}

A rotating LiDAR acquires points at different robot poses within one scan,
causing geometric inconsistency. Deskewing corrects each point to a common
timestamp using IMU-interpolated poses.
IMU preintegration~\cite{forster2017manifold} accumulates measurements between
two keyframes into a single relative rotation on $SO(3)$:

\begin{equation}
  \Delta\hat{\mathbf{R}}_{ij} =
    \prod_{k=i}^{j-1} \operatorname{Exp}\!\bigl(\tilde{\boldsymbol{\omega}}_k\,\Delta t\bigr),
  \label{eq:imu_preint}
\end{equation}

\noindent where $\operatorname{Exp}: \mathfrak{so}(3) \to SO(3)$ is the matrix
exponential. This constraint is fixed across re-optimisation iterations.
FAST-LIO2 and LIO-SAM both rely on deskewing and preintegration.

% ─────────────────────────────────────────────────────────────────────────────
\section{Visual SLAM Fundamentals}

Visual SLAM estimates trajectory from image sequences by tracking sparse
features. This thesis evaluates feature-based stereo systems that recover
depth from stereo disparity (Section~\ref{sec:sensors}).

\subsection{Stereo Vision and Triangulation}

Stereo vision~\cite{siegwart2004introduction} recovers the 3D position of a
scene point from its disparity using Equation~\eqref{eq:stereo_depth}.
Each matched feature pair is lifted to a 3D landmark by back-projection.
Depth uncertainty grows quadratically with distance, so distant landmarks are
typically discarded.

\subsection{Feature Detection and Description}

Visual keypoint detectors identify salient regions and assign binary descriptors
for matching.
GFTT selects corners by maximising the smaller eigenvalue of the local gradient
matrix and is used in RTABMap's front-end.
ORB~\cite{rublee2011orb} combines FAST keypoints with a rotation-steered BRIEF
descriptor, giving rotation and scale invariance; it is used by ORB-SLAM3 and
RTABMap.

\subsection{Visual Odometry}

Visual odometry~\cite{scaramuzza2011visual} estimates incremental motion from
matched feature tracks via PnP or essential matrix decomposition.
Local bundle adjustment refines recent poses by minimising reprojection error:

\begin{equation}
  E = \sum_{i,j} \rho\!\left(
        \bigl\lVert \mathbf{u}_{ij} - \pi(\mathbf{T}_i,\, \mathbf{l}_j) \bigr\rVert^2
      \right),
  \label{eq:reprojection}
\end{equation}

\noindent where $\mathbf{u}_{ij}$ is the observed keypoint, $\pi$ projects
landmark $\mathbf{l}_j$ through pose $\mathbf{T}_i \in SE(3)$, and $\rho$ is
a robust cost function. Errors accumulate without loop closure.

\subsection{Bag-of-Words Place Recognition}
\label{sec:bow}

BoW place recognition~\cite{galvez2012bags} represents each image as a sparse
histogram over a visual vocabulary built offline by clustering descriptors.
At runtime the current histogram is compared to stored ones by cosine similarity;
high-scoring candidates are verified geometrically with RANSAC before a
loop-closure edge is added.
RTABMap uses DBoW3; ORB-SLAM3 uses DBoW2.

% ─────────────────────────────────────────────────────────────────────────────
\section{Evaluation Metrics}
\label{sec:eval_metrics}

Trajectory metrics quantify how closely the estimated path matches simulator
ground-truth poses; map metrics quantify the geometric fidelity of the
reconstructed 3D point cloud.

\subsection{Umeyama Alignment}

Before computing errors, trajectories are brought into the same frame by
Umeyama alignment~\cite{umeyama1991least}, which minimises:

\begin{equation}
  \varepsilon^2(\mathbf{R}, \mathbf{t}, c) =
    \frac{1}{n}\sum_{i=1}^{n}
    \lVert \mathbf{q}_i - (c\mathbf{R}\mathbf{p}_i + \mathbf{t}) \rVert^2,
  \label{eq:umeyama_cost}
\end{equation}

\noindent where $\mathbf{p}_i$ are estimated positions and $\mathbf{q}_i$ are
ground-truth positions. For metric sensors the scale is fixed at $c = 1$,
reducing alignment to $SE(3)$.

\subsection{Absolute Trajectory Error (ATE)}

ATE~\cite{sturm2012benchmark} measures global consistency as the RMSE of
per-pose translational errors after alignment:

\begin{equation}
  \text{ATE} = \sqrt{\frac{1}{n}\sum_{i=1}^{n}
               \lVert \hat{\mathbf{p}}_i - \mathbf{p}_i^* \rVert^2}.
  \label{eq:ate}
\end{equation}

\noindent A large ATE indicates drift or absent loop closure. Mean and maximum
errors are also reported.

\subsection{Relative Pose Error (RPE)}

RPE~\cite{sturm2012benchmark} measures local odometric accuracy by comparing
consecutive relative motions between estimate and ground truth:

\begin{equation}
  \text{RPE} = \sqrt{\frac{1}{n-1}\sum_{i=1}^{n-1}
               \lVert \mathbf{e}_{i}^{\,\text{trans}} \rVert^2}.
  \label{eq:rpe}
\end{equation}

\noindent RPE captures short-term drift that ATE may hide after loop closure.

\subsection{Chamfer Distance}

Chamfer Distance~\cite{fan2017point} measures geometric fidelity between the
estimated map $\mathcal{A}$ and reference $\mathcal{B}$:

\begin{equation}
  \text{CD}(\mathcal{A},\mathcal{B}) =
    \frac{1}{|\mathcal{A}|}\sum_{\mathbf{a}\in\mathcal{A}}
      \min_{\mathbf{b}}\lVert\mathbf{a}-\mathbf{b}\rVert
    +
    \frac{1}{|\mathcal{B}|}\sum_{\mathbf{b}\in\mathcal{B}}
      \min_{\mathbf{a}}\lVert\mathbf{b}-\mathbf{a}\rVert.
  \label{eq:chamfer}
\end{equation}

\noindent CD is reported in metres but is sensitive to outlier points.

\subsection{F-score}

The F-score~\cite{knapitsch2017tanks} combines precision and recall at a
distance threshold $d$:

\begin{align}
  P(d) &= \frac{1}{|\mathcal{A}|}
          \bigl|\{\mathbf{a}\in\mathcal{A} :
          \min_{\mathbf{b}}\lVert\mathbf{a}-\mathbf{b}\rVert < d\}\bigr|, \\[4pt]
  R(d) &= \frac{1}{|\mathcal{B}|}
          \bigl|\{\mathbf{b}\in\mathcal{B} :
          \min_{\mathbf{a}}\lVert\mathbf{b}-\mathbf{a}\rVert < d\}\bigr|, \\[4pt]
  F(d) &= \frac{2\,P(d)\,R(d)}{P(d)+R(d)}.
  \label{eq:fscore}
\end{align}

Three thresholds are used: $d = 5$\,cm (strict), $d = 10$\,cm (moderate), and
$d = 20$\,cm (lenient). Unlike CD, the F-score is bounded to $[0,1]$ and is not
dominated by outliers.
